\documentclass[12pt]{article}
\usepackage[margin=0.75in]{geometry}
\usepackage{lmodern}
\usepackage{microtype}
\usepackage{multicol}
\usepackage{fancyhdr}
\usepackage{titlesec}
\usepackage{enumitem}
\usepackage[hidelinks]{hyperref}
\usepackage{parskip}

\setlength{\columnsep}{0.28in}
\setlength{\parindent}{1.05em}
\setlength{\parskip}{0.35em}
\setlength{\headheight}{15pt}
\titleformat{\section}{\large\bfseries\scshape}{}{0pt}{}
\titleformat{\subsection}{\normalsize\bfseries}{}{0pt}{}
\pagestyle{fancy}
\fancyhf{}
\fancyhead[L]{Newman: Knowledge Its Own End}
\fancyhead[R]{Reading Handout}
\fancyfoot[C]{\thepage}
\renewcommand{\headrulewidth}{0.4pt}

\begin{document}
\begin{center}
{\LARGE\bfseries Knowledge and the Free Mind}\par\vspace{0.35em}
{\large\scshape John Henry Newman on Liberal Education}\par\vspace{0.6em}
\rule{0.82\textwidth}{0.4pt}\par\vspace{0.9em}
\end{center}

\section*{Vocabulary Preview}
\begin{multicols}{2}
\begin{itemize}[leftmargin=*, itemsep=0.2em]
\item \textbf{Liberal education}: education that frees and forms the mind, not merely job training.
\item \textbf{Utility}: practical usefulness; what something can be used to produce.
\item \textbf{Intellect}: the mind's power to understand, judge, reason, and connect ideas.
\item \textbf{Cultivation}: careful development or training toward excellence.
\item \textbf{Philosophical habit}: a trained habit of seeing causes, relations, distinctions, and principles.
\item \textbf{Enlargement}: expansion; making the mind broader and more capable.
\item \textbf{Discipline}: training that strengthens a habit or faculty.
\item \textbf{Professional}: connected to a particular job, career, or trade.
\item \textbf{Servile}: fit for a servant; narrow, dependent, or lacking freedom.
\item \textbf{End}: a goal, purpose, or final object sought for its own sake.
\item \textbf{Equitable}: fair-minded; able to judge without narrow prejudice.
\item \textbf{Ulterior}: beyond what is immediate; further or hidden.
\end{itemize}
\end{multicols}

\section*{Introduction}
John Henry Newman lived from 1801 to 1890. In \textit{The Idea of a University}, he asks what education is for. Is school mainly for practical training, careers, money, and useful skills? Or is there a kind of knowledge that is valuable because it forms the mind itself?

Newman does not despise useful knowledge. He knows that medicine, law, engineering, business, and trades matter. But he argues that a human being needs more than training for a function. A free person needs a cultivated intellect: a mind able to see relationships, weigh evidence, judge fairly, and understand truth beyond immediate usefulness. This final reading brings the unit back to Plato's cave: education should not merely hand students tools; it should help turn the mind toward reality.

\section*{Source Note}
Adapted and modernized from John Henry Newman, \textit{The Idea of a University}, Discourse V, ``Knowledge Its Own End.'' This handout uses focused passages on liberal knowledge, utility, intellectual cultivation, and the philosophical habit of mind. The language has been modernized for student reading while preserving Newman's sequence, claims, and central images.

\begin{multicols}{2}
\section*{Reading: Knowledge Its Own End}

A university may be considered in two ways: by looking at its studies, or by looking at its students. I have already argued that all knowledge forms one whole, and that the separate sciences are parts of that whole. Now I turn to the students. What kind of education does a university give them? And in what sense can such education be called useful?

There is no branch of knowledge that looks exactly the same when it is seen by itself as when it is seen as part of the whole. One science completes, corrects, and balances another. Each subject has its own light, but each also casts its own shadow when it is isolated from the rest. If a student learns one field without any sense of its relation to others, that field may become narrow, exaggerated, or misleading.

For this reason, it matters that a university contains many kinds of study. No student can pursue every subject deeply. Still, students gain something by living among teachers and learners who represent the whole circle of knowledge. Scholars who love different fields must learn to respect one another, consult one another, and adjust the claims of their subjects in relation to other truths.

From this common life of learning, a clear atmosphere of thought is created. The student breathes that atmosphere even if he studies only a few subjects. He begins to grasp the great outlines of knowledge: its principles, its parts, its greater and lesser matters, its lights and shades. He learns not merely facts, but proportion and judgment.

This is why such education is called liberal. It forms a habit of mind that lasts through life. Its qualities are freedom, fairness, calmness, moderation, and wisdom. It gives what I have elsewhere called a philosophical habit. This is the special fruit of a university education.

Now practical people will ask: What is the gain of this philosophical habit? What does it do? Where does it end? How does it profit? Particular sciences lead to definite arts. Medicine leads to healing; law leads to legal practice; engineering leads to construction; agriculture leads to crops. But what is the art produced by this science of sciences? What practical fruit comes from philosophy in this broad sense?

My answer is that liberal knowledge has a real and sufficient end, but that end cannot be separated from the knowledge itself. Knowledge is capable of being its own end. The human mind is made in such a way that any real knowledge is, in some sense, its own reward. This is especially true of philosophical knowledge: knowledge that sees truth in its branches, relations, mutual bearings, and proper values.

I am not claiming here that such knowledge is worth more than wealth, power, honor, comfort, or the conveniences of life. I am saying that knowledge is a real good in its own nature. It is worth thought and labor because it perfects something in us.

Some people assume that knowledge must be pursued for something beyond itself. They say it must be useful either for this world or for the next. If it serves worldly needs, they call it useful knowledge. If it serves the soul directly, they call it religious knowledge. If liberal knowledge does neither of these directly, they conclude that it is not good at all.

But that conclusion is mistaken. Useful knowledge has done great work, and I do not deny it. It has increased health, comfort, travel, trade, and ordinary well-being. We owe much to practical discoveries. But liberal knowledge should not be judged by the same standard, as if its direct business were to build machines, manage markets, or train a person for one profession.

Nor should we pretend that liberal knowledge, by itself, makes a person virtuous or holy. Knowledge is one thing; virtue is another. Good sense is not the same as conscience. Refinement is not humility. A wide and just view of things is not the same as faith. Philosophy can enlighten the mind without conquering the passions. A cultivated intellect is good, but it is not the whole of goodness.

So what is liberal education in itself? It is the cultivation of the intellect as intellect. Its object is intellectual excellence. Every kind of thing has its own perfection. A garden is not cultivated only to become an orchard or a field of grain. We care for walks, trees, water, lawns, and slopes because there is a special beauty in their order as a whole. A city, a palace, a church, or a poem may have a beauty that leads to nothing beyond itself.

In the same way, the intellect has its own beauty and perfection. To open the mind, to correct it, to refine it, to teach it to know, digest, master, rule, and use its knowledge: this is an intelligible object. A trained mind gains command over its own powers. It becomes able to apply itself, to change direction when needed, to use method, to judge accurately, to see quickly, to find resources, and to express itself well.

This does not make education useless. It gives education a deeper use than immediate usefulness. A person may be trained for a job and still be narrow, easily misled, or unable to judge the larger meaning of what he does. A liberal education forms the mind that uses all other knowledge. It helps a person see more than fragments. It gives freedom from the tyranny of slogans, fashions, and isolated opinions.

Therefore knowledge has dignity beyond profit. It is not merely a servant of money, power, or career. Liberal education seeks the perfection of the mind. It prepares a person to see relations, weigh claims, resist confusion, and pursue truth with steadiness. If shadows come from appearances, idols, disordered loves, and unjust power, then education must do more than provide information. It must form a free mind capable of seeking reality.

\section*{Reading Focus}
As you reread, mark the text:
\begin{enumerate}[leftmargin=*]
\item \textbf{Underline} one sentence that defines liberal education.
\item \textbf{Circle} one sentence that answers the ``What is the use?'' objection.
\item \textbf{Put a star} beside one sentence that connects education with freedom of mind.
\item \textbf{Write ``whole''} beside a place where Newman says knowledge must be seen in relation to other knowledge.
\end{enumerate}

\section*{Utility vs. Liberal Education}
Write two short speeches, three to five sentences each.
\begin{enumerate}[leftmargin=*]
\item \textbf{Speaker 1: Education should be useful.} Argue that school should mainly prepare students for jobs, practical skills, money, and ordinary life.
\item \textbf{Speaker 2: Education should be liberal.} Argue that school should form judgment, breadth of mind, and love of truth beyond immediate usefulness.
\end{enumerate}
Make each speaker strong. Do not caricature either side.

\section*{Unit Synthesis}
In your notebook, complete the following statements, then write one sentence explaining how all five readings fit together.
\begin{itemize}[leftmargin=*]
\item Plato says education is like \underline{\hspace{0.7in}}.
\item Bacon says our minds are distorted by \underline{\hspace{0.7in}}.
\item Augustine says we resist truth because \underline{\hspace{0.7in}}.
\item Douglass says knowledge is dangerous to unjust power because \underline{\hspace{0.7in}}.
\item Newman says education should form \underline{\hspace{0.7in}}.
\end{itemize}

\section*{Final Claim Preparation}
Choose one claim you could defend in a paper. List two readings that support it, one reading that challenges or complicates it, and one passage you would use as evidence.
\begin{enumerate}[leftmargin=*]
\item People resist truth mainly because truth is painful.
\item People resist truth mainly because their minds are distorted by idols.
\item People resist truth mainly because they love the wrong things.
\item People resist truth mainly because someone benefits from keeping them ignorant.
\item Education frees people from shadows only if it forms the whole mind and soul.
\end{enumerate}

\section*{Handout Summary}
End by writing a short summary of this handout. In three to five sentences, explain Newman's main claim about liberal education and how it connects to the unit question: Why do people sometimes love shadows more than truth?

\end{multicols}
\end{document}
